\documentclass[../../main.tex]{subfiles}
\graphicspath{{../../image/}{../../figure/}}
\begin{document}

\section{稀疏波的散开}\label{section:ch12-稀疏波的散开}

在原书第 11.2 节中我们看到，稀疏波的散开会对真非线性标量守恒律的解诱导单侧 Lipschitz 条件。
此处，在两守恒律组\eqref{align:ch12-系统}的语境下，我们将遇到类似效应。我们将看到，1-（或 2-）稀疏波的散开起到降低相应 Riemann 不变量 $z$（或 $w$）的下降（或上升）斜率的作用。由于介入的波相互作用，这一机制不再能像标量情形那样维持单侧 Lipschitz 条件；但它仍能独立于初始数据保持解的总变差有界。

我们再次考虑上一节中讨论过的、具有小振荡\eqref{align:ch12-解小振荡}的解 $(z,w)$。主要结果是

\begin{theorem}\label{theorem:ch12-TV散开}
对任意 $-\infty<x<y<\infty$ 及 $t>0$，
\begin{align}
TV_{[x,y]}z(\cdot,t)+TV_{[x,y]}w(\cdot,t)\leqslant b\frac{y-x}{t}+\beta\delta,
\label{align:ch12-TV散开}
\end{align}
其中 $b$ 与 $\beta$ 是可依赖于 $F$ 但独立于初始数据的常数。
\end{theorem}

上述定理的证明将分成若干步。\cref{section:ch12-解的界}中引入的记号在此自由使用。特别地，如前，$c$ 表示可依赖于 $F$ 但独立于 $\delta$ 的一般常数。

\begin{lemma}\label{lemma:ch12-散开估计}
固定 $\bar t>0$，任取 $-\infty<x_\ell<x_r<\infty$，使 $x_r-x_\ell$ 相对于 $\bar t$ 充分小。构造从 $(x_\ell,\bar t)$、$(x_r,\bar t)$ 发出的极小（或极大）后向 1-（或 2-）特征线 $\xi_\ell(\cdot)$、$\xi_r(\cdot)$（或 $\zeta_\ell(\cdot)$、$\zeta_r(\cdot)$），令 $\mathscr{F}$（或 $\mathscr{G}$）表示由 $\xi_\ell$、$\xi_r$（或 $\zeta_\ell$、$\zeta_r$）与时间线 $t=\bar t$、$t=\bar t/2$ 的图像所围成的区域。则
\begin{align}
z(x_\ell,\bar t)-z(x_r,\bar t)&\leqslant\hat c\exp(\bar c\delta\bar V)\frac{x_r-x_\ell}{\bar t}+\mathscr{P}_2(\mathscr{F})+\mathscr{Q}_2(\mathscr{F}),
\label{align:ch12-散开估计1}\\
w(x_r,\bar t)-w(x_\ell,\bar t)&\leqslant\hat c\exp(\bar c\delta\bar V)\frac{x_r-x_\ell}{\bar t}+\mathscr{P}_1(\mathscr{G})+\mathscr{Q}_1(\mathscr{G}),
\label{align:ch12-散开估计2}
\end{align}
其中 $\bar V$ 表示 $w$（或 $z$）沿 $\xi_\ell(\cdot)$（或 $\zeta_r(\cdot)$）在区间 $[\frac12\bar t,\bar t]$ 上的迹的总变差。
\end{lemma}

\begin{proof}
只需证明\eqref{align:ch12-散开估计1}。令 $(z_\ell(\cdot),w_\ell(\cdot))$ 与 $(z_r(\cdot),w_r(\cdot))$ 分别表示沿 $\xi_\ell(\cdot)$ 与 $\xi_r(\cdot)$ 的 $(z,w)$ 迹。

我们考虑解值域上特征速度 $\mu(z,w)$ 的下确界 $\tilde\mu$ 与上确界 $\bar\mu$。从点 $(\xi_r(t),t)$（$t\in[\frac12\bar t,\bar t]$）发出的斜率为 $\tilde\mu$ 与 $\bar\mu$ 的直线分别在时刻 $f(t)$ 与 $g(t)$ 被 $\xi_\ell(\cdot)$ 截取。函数 $f$ 与 $g$ 均为斜率 $1+O(\delta)$ 的 Lipschitz 函数，且
\begin{align}
0\leqslant g(t)-f(t)\leqslant c_1\delta[\xi_r(t)-\xi_\ell(f(t))].
\label{align:ch12-fg}
\end{align}
映射将 $(\xi_r(t),t)$ 映到 $(\xi_\ell(f(t)),f(t))$，诱导 $\xi_\ell$ 与 $\xi_r$ 图像上点的配对。由
\begin{align}
\xi_r(t)-\xi_\ell(f(t))=\tilde\mu[t-f(t)],
\label{align:ch12-配对}
\end{align}
得
\begin{align}
\dot f(t)&=1-\frac{1}{\tilde\mu-\dot\xi_\ell(f(t))}[\dot\xi_r(t)-\dot\xi_\ell(f(t))],
\label{align:ch12-fdot}\\
\frac{d}{dt}[\xi_r(t)-\xi_\ell(f(t))]&=\frac{\tilde\mu}{\tilde\mu-\dot\xi_\ell(f(t))}[\dot\xi_r(t)-\dot\xi_\ell(f(t))],
\label{align:ch12-dfdot}
\end{align}
在 $[\frac12\bar t,\bar t]$ 上几乎处处成立。为从下方估计\eqref{align:ch12-dfdot}右端，我们从
\begin{align}
\dot\xi_r(t)-\dot\xi_\ell(f(t))&=\lambda(z_r(t),w_r(t))-\lambda(z_\ell(f(t)),w_\ell(f(t)))\nonumber\\
&=\bar\lambda_z[z_r(t)-z_\ell(f(t))]+\bar\lambda_w[w_r(t)-w_\ell(f(t))].
\label{align:ch12-速度差}
\end{align}
出发。由\cref{theorem:ch12-传播主定理}，
\begin{align}
z_r(t)-z_\ell(f(t))\leqslant z(x_r,\bar t)-z(x_\ell,\bar t)-\sum[z_r(\tau+)-z_r(\tau-)],
\label{align:ch12-z差-s06}
\end{align}
其中求和遍历区间 $(t,\bar t)$ 内 $z_r(\cdot)$ 的跳跃点集。如\cref{lemma:ch12-迹差估计}的证明，对每个此类跳跃点 $\tau$，可联系由单节点 $(\xi_r(\tau),\tau)$ 组成的平凡 2-特征线树 $\mathscr{N}$，从而推得
\begin{align}
-\sum[z_r(\tau+)-z_r(\tau-)]\leqslant\mathscr{P}_2(\mathscr{F})+\mathscr{Q}_2(\mathscr{F}).
\label{align:ch12-z跳跃}
\end{align}
对 $t\in[\frac12\bar t,\bar t]$，我们构造从 $(\xi_r(t),t)$ 发出的极大后向 2-特征线，它在时刻 $h(t)$ 被 $\xi_\ell(\cdot)$ 截取；$f(t)\leqslant h(t)\leqslant g(t)$。由\cref{theorem:ch12-传播主定理}，$w_\ell(h(t))\geqslant w_r(t)$，故
\begin{align}
w_r(t)-w_\ell(f(t))\leqslant w_\ell(h(t))-w_\ell(f(t))\leqslant V(f(t))-V(g(t)),
\label{align:ch12-w差}
\end{align}
其中 $V(\tau)$ 度量 $w_\ell(\cdot)$ 在区间 $[\tau,\bar t]$ 上的总变差。

现将\eqref{align:ch12-dfdot}在区间 $(s,\bar t)$ 上积分。回忆 $\bar\lambda_z<0$、$\bar\lambda_w<0$，结合\eqref{align:ch12-速度差}、\eqref{align:ch12-z差}、\eqref{align:ch12-z跳跃}与\eqref{align:ch12-w差}，推得
\begin{align}
\xi_r(s)-\xi_\ell(f(s))&\leqslant\xi_r(\bar t)-\xi_\ell(f(\bar t))\nonumber\\
&\quad+c_2^{-1}(\bar t-s)[z(x_r,\bar t)-z(x_\ell,\bar t)+\mathscr{P}_2(\mathscr{F})+\mathscr{Q}_2(\mathscr{F})]\nonumber\\
&\quad+c_3\int_s^{\bar t}[V(f(t))-V(g(t))]\,dt.
\label{align:ch12-积分}
\end{align}
交换积分次序，
\begin{align}
\int_s^{\bar t}[V(f(t))-V(g(t))]\,dt&=-\int_s^{\bar t}\int_{f(t)}^{g(t)}dV(\tau)\,dt\nonumber\\
&\leqslant-\int_{f(s)}^{f(\bar t)}[f^{-1}(\tau)-g^{-1}(\tau)]dV(\tau)-\int_{f(\bar t)}^{g(\bar t)}[\bar t-g^{-1}(\tau)]dV(\tau)\nonumber\\
&=-\int_s^{\bar t}[t-g^{-1}(f(t))]dV(f(t))-\int_{f(\bar t)}^{g(\bar t)}[\bar t-g^{-1}(\tau)]dV(\tau).
\label{align:ch12-交换}
\end{align}
由\eqref{align:ch12-fg}，
\begin{align}
t-g^{-1}(f(t))\leqslant c_4\delta[\xi_r(t)-\xi_\ell(f(t))],\qquad \frac{\bar t}{2}\leqslant t\leqslant\bar t,
\label{align:ch12-g1}\\
\bar t-g^{-1}(\tau)\leqslant c_4\delta[\xi_r(\bar t)-\xi_\ell(f(\bar t))],\qquad f(\bar t)\leqslant\tau\leqslant g(\bar t),
\label{align:ch12-g2}
\end{align}
故\eqref{align:ch12-积分}给出
\begin{align}
\xi_r(s)-\xi_\ell(f(s))&\leqslant\exp(c_3c_4\delta\bar V)[\xi_r(\bar t)-\xi_\ell(f(\bar t))]\nonumber\\
&\quad+c_2^{-1}(\bar t-s)[z(x_r,\bar t)-z(x_\ell,\bar t)+\mathscr{P}_2(\mathscr{F})+\mathscr{Q}_2(\mathscr{F})]\nonumber\\
&\quad-c_3c_4\delta\int_s^{\bar t}[\xi_r(t)-\xi_\ell(f(t))]dV(f(t)),
\label{align:ch12-Gronwall前}
\end{align}
对任意 $s\in[\frac34\bar t,\bar t]$ 成立。积分上述 Gronwall 型不等式，得
\begin{align}
\xi_r(s)-\xi_\ell(f(s))&\leqslant\exp(2c_3c_4\delta\bar V)[\xi_r(\bar t)-\xi_\ell(f(\bar t))]\nonumber\\
&\quad+c_2^{-1}\int_s^{\bar t}\exp\{c_3c_4\delta[V(f(s))-V(f(\tau))]\}\,d\tau\,[z(x_r,\bar t)-z(x_\ell,\bar t)+\mathscr{P}_2(\mathscr{F})+\mathscr{Q}_2(\mathscr{F})].
\label{align:ch12-Gronwall}
\end{align}
我们取 $s=\frac34\bar t$ 应用\eqref{align:ch12-Gronwall}。\eqref{align:ch12-Gronwall}左端非负，且 $\xi_r(\bar t)-\xi_\ell(f(\bar t))\leqslant c_5(x_r-x_\ell)$。因此\eqref{align:ch12-Gronwall}蕴含\eqref{align:ch12-散开估计1}，常数 $\bar c=2c_3c_4$、$\hat c=4c_2c_5$。证明完成。
\end{proof}

下文我们将在假设\eqref{align:ch12-散开估计1}、\eqref{align:ch12-散开估计2}中出现的常数 $\bar V$ 满足
\begin{align}
\bar c\delta\bar V\leqslant\log 2
\label{align:ch12-barV}
\end{align}
之下进行操作。由\cref{lemma:ch12-相互作用有界}，当初始数据以充分小的 $a$ 满足\eqref{align:ch12-初值总变差}时，此式必然成立。进而，由于有限依赖域性质，即使当初始数据仅有局部有界变差并以充分小的 $\delta$ 满足\eqref{align:ch12-初值上确界}时，\eqref{align:ch12-barV} 对充分小的 $\bar t$ 仍成立。下文将证明，只要初始数据的振荡充分小（即 $\delta$ 小），\eqref{align:ch12-barV} 实际上对任意 $\bar t>0$ 成立。

\begin{lemma}\label{lemma:ch12-NVPV}
对任意 $-\infty<\bar x<\bar y<\infty$ 及 $\bar t>0$，
\begin{align}
NV_{[\bar x,\bar y]}z(\cdot,\bar t)+PV_{[\bar x,\bar y]}w(\cdot,\bar t)&\leqslant 4\hat c\frac{\bar y-\bar x}{\bar t}\nonumber\\
&\quad+c\delta^2\left\{TV_{[\bar x-\frac12\bar\mu\bar t,\bar y-\frac12\bar\lambda\bar t]}z\left(\cdot,\frac{\bar t}{2}\right)+TV_{[\bar x-\frac12\bar\mu\bar t,\bar y-\frac12\bar\lambda\bar t]}w\left(\cdot,\frac{\bar t}{2}\right)\right\},
\label{align:ch12-NVPV1}\\
TV_{[\bar x,\bar y]}z(\cdot,\bar t)+TV_{[\bar x,\bar y]}w(\cdot,\bar t)&\leqslant 8\hat c\frac{\bar y-\bar x}{\bar t}+8\delta\nonumber\\
&\quad+c\delta^2\left\{TV_{[\bar x-\frac12\bar\mu\bar t,\bar y-\frac12\bar\lambda\bar t]}z\left(\cdot,\frac{\bar t}{2}\right)+TV_{[\bar x-\frac12\bar\mu\bar t,\bar y-\frac12\bar\lambda\bar t]}w\left(\cdot,\frac{\bar t}{2}\right)\right\},
\label{align:ch12-NVPV2}
\end{align}
其中 $\bar\lambda$ 是解值域上 $\lambda(z,w)$ 的下确界，$\bar\mu$ 是 $\mu(z,w)$ 的上确界。
\end{lemma}

\begin{proof}
结合\eqref{align:ch12-散开估计1}、\eqref{align:ch12-散开估计2}、\eqref{align:ch12-barV}与\cref{lemma:ch12-不相交求和}，立即推出
\begin{align}
NV_{[\bar x,\bar y]}z(\cdot,\bar t)+PV_{[\bar x,\bar y]}w(\cdot,\bar t)\leqslant 4\hat c\frac{\bar y-\bar x}{\bar t}+2[\mathscr{P}_1(\mathscr{H})+\mathscr{Q}_1(\mathscr{H})+\mathscr{P}_2(\mathscr{H})+\mathscr{Q}_2(\mathscr{H})],
\label{align:ch12-NVPV前}
\end{align}
其中 $\mathscr{H}$ 表示由从 $(\bar y,\bar t)$ 发出的极小后向 1-特征线 $\xi(\cdot)$ 的图像、从 $(\bar x,\bar t)$ 发出的极大后向 2-特征线 $\zeta(\cdot)$ 的图像，以及时间线 $t=\bar t$、$t=\bar t/2$ 所围成的区域。

我们通过将\cref{lemma:ch12-相互作用有界}应用于时间原点由 $t=0$ 平移到 $t=\bar t/2$ 的情形来估计 $\mathscr{P}_1(\mathscr{H})+\mathscr{Q}_1(\mathscr{H})+\mathscr{P}_2(\mathscr{H})+\mathscr{Q}_2(\mathscr{H})$。这给出\eqref{align:ch12-NVPV1}。

由于总变差是负变差与正变差之和，而负变差与正变差之差被振荡控制，\eqref{align:ch12-NVPV1} 连同\eqref{align:ch12-解小振荡}给出\eqref{align:ch12-NVPV2}。证明完成。
\end{proof}

\begin{proof}[\cref{theorem:ch12-TV散开} 的证明]
为建立\eqref{align:ch12-TV散开}，先取 $\bar t=t$、$\bar x=x$、$\bar y=y$ 写入\eqref{align:ch12-NVPV2}。为估计所得不等式右端，我们对 $\bar t=\frac12t$、$\bar x=x-\frac12\bar\mu t$、$\bar y=y-\frac12\bar\lambda t$ 再次应用\eqref{align:ch12-NVPV2}。这给出
\begin{align}
TV_{[x-\frac12\bar\mu t,y-\frac12\bar\lambda t]}z\left(\cdot,\frac{t}{2}\right)+TV_{[x-\frac12\bar\mu t,y-\frac12\bar\lambda t]}w\left(\cdot,\frac{t}{2}\right)
&\leqslant 16\hat c\frac{y-x}{t}+8\hat c(\bar\mu-\bar\lambda)+8\delta\nonumber\\
&\quad+c\delta^2\left\{TV_{[x-\frac34\bar\mu t,y-\frac34\bar\lambda t]}z\left(\cdot,\frac{t}{4}\right)+TV_{[x-\frac34\bar\mu t,y-\frac34\bar\lambda t]}w\left(\cdot,\frac{t}{4}\right)\right\}.
\label{align:ch12-迭代1}
\end{align}
类似地，为估计\eqref{align:ch12-迭代1}右端，我们对 $\bar t=\frac14t$、$\bar x=x-\frac34\bar\mu t$、$\bar y=y-\frac34\bar\lambda t$ 应用\eqref{align:ch12-NVPV2}。于是得
\begin{align}
TV_{[x-\frac34\bar\mu t,y-\frac34\bar\lambda t]}z\left(\cdot,\frac{t}{4}\right)+TV_{[x-\frac34\bar\mu t,y-\frac34\bar\lambda t]}w\left(\cdot,\frac{t}{4}\right)
&\leqslant 32\hat c\frac{y-x}{t}+24\hat c(\bar\mu-\bar\lambda)+8\delta\nonumber\\
&\quad+c\delta^2\left\{TV_{[x-\frac78\bar\mu t,y-\frac78\bar\lambda t]}z\left(\cdot,\frac{t}{8}\right)+TV_{[x-\frac78\bar\mu t,y-\frac78\bar\lambda t]}w\left(\cdot,\frac{t}{8}\right)\right\}.
\label{align:ch12-迭代2}
\end{align}
继续下去并取极限，即得\eqref{align:ch12-TV散开}，其中
\begin{align}
b=\frac{8\hat c}{1-2c\delta^2},\qquad \beta=\frac{8}{1-c\delta^2}+\frac{8c\hat c\delta(\bar\mu-\bar\lambda)}{(1-c\delta^2)(1-2c\delta^2)}.
\label{align:ch12-bbeta}
\end{align}

上述推导以\eqref{align:ch12-barV} 成立为前提；因此为完成证明，现须验证此条件。回忆\cref{lemma:ch12-散开估计}中 $\bar V$ 的定义，并将\cref{lemma:ch12-相互作用有界}用于时间原点由 $0$ 平移到 $\frac12t$ 的情形，推得
\begin{align}
\bar V\leqslant c\sup_{\bar x}\left\{TV_{[\bar x-\frac12\bar\mu t,\bar x-\frac12\bar\lambda t]}z\left(\cdot,\frac{t}{2}\right)+TV_{[\bar x-\frac12\bar\mu t,\bar x-\frac12\bar\lambda t]}w\left(\cdot,\frac{t}{2}\right)\right\}.
\label{align:ch12-barV估计}
\end{align}
我们以\eqref{align:ch12-TV散开}估计\eqref{align:ch12-barV估计}右端，得
\begin{align}
\bar V\leqslant cb(\bar\mu-\bar\lambda)+c\beta\delta,
\label{align:ch12-barV终}
\end{align}
故当 $\delta$ 充分小时\eqref{align:ch12-barV} 确实满足。证明完成。
\end{proof}

我们现证明：振荡充分小但总变差任意大的初始数据，诱导出本节全程假设的 $L^\infty$ 界\eqref{align:ch12-解小振荡}。

\begin{theorem}\label{theorem:ch12-小振荡Linf}
存在仅依赖于 $F$ 的正常数 $\gamma$，使得由小振荡
\begin{align}
|z(x,0)|+|w(x,0)|<\gamma\delta^2,\qquad -\infty<x<\infty,
\label{align:ch12-小振荡}
\end{align}
但总变差不受限的初始数据生成的解满足\eqref{align:ch12-解小振荡}。
\end{theorem}

\begin{proof}
假设\eqref{align:ch12-小振荡} 成立且 $\gamma$ 充分小，我们将证明在上半平面 $-\delta<z(x,t)<\delta$ 与 $-\delta<w(x,t)<\delta$。反证，设上述四个不等式之一在某点被破坏，例如 $z(\bar x,\bar t)\geqslant\delta$。

我们由 $8\hat c(\bar y-\bar x)=\delta\bar t$ 确定 $\bar y$，其中 $\hat c$ 是\eqref{align:ch12-散开估计1}中出现的常数，并应用\eqref{align:ch12-NVPV1}。\eqref{align:ch12-NVPV1}右端第一项此处被 $\frac14\delta$ 界定；第二项由\cref{theorem:ch12-TV散开}被 $\tilde c\delta^2$ 界定。因此，对充分小的 $\delta$，$z(\cdot,\bar t)$ 在区间 $[\bar x,\bar y]$ 上的负（下降）变差不超过 $\frac12\delta$。由此，对所有 $x\in[\bar x,\bar y]$，$z(x,\bar t)\geqslant\frac12\delta$。特别地，
\begin{align}
\int_{\bar x}^{\bar y}[|z(x,\bar t)|+|w(x,\bar t)|]\,dx\geqslant(\bar y-\bar x)\frac{\delta}{2}=\frac{1}{16\hat c}\delta^2\bar t.
\label{align:ch12-L1下界}
\end{align}
现诉诸 $L^1$ 估计\eqref{align:ch12-L1}（将在\cref{section:ch12-L1估计}中由引理 12.8.2 建立），
并结合\eqref{align:ch12-小振荡}推得
\begin{align}
\int_{\bar x}^{\bar y}[|z(x,\bar t)|+|w(x,\bar t)|]\,dx\leqslant 4[(\bar y-\bar x)+2c\bar t]\gamma\delta^2=\gamma\left[\frac{\delta}{2\hat c}+8c\right]\delta^2\bar t.
\label{align:ch12-L1上界}
\end{align}
显然，当 $\gamma$ 充分小时，\eqref{align:ch12-L1下界} 与\eqref{align:ch12-L1上界} 不相容，这给出期望的矛盾。证明完成。
\end{proof}

结合原书第 1.7 节所述 $BV$ 函数的紧性性质，
估计\eqref{align:ch12-TV散开}表明：从具有局部有界变差初始数据的解出发，可通过完备化在仅为 $L^\infty$、振荡充分小的初始数据下构造 $BV_{\mathrm{loc}}$ 解。因此，真非线性两守恒律组的解算子通过已在真非线性标量守恒律（定理 11.2.2）语境中遇到的机制对初始数据进行正则化。

\end{document}
